\section{Introduction}\label{sec:intro}

In Bernoulli bond percolation with parameter $p\in[0,1]$ on an infinite connected graph, every edge is
declared open with probability~$p$, independently of all other edges, and one studies the connected
components of the open subgraph, the \emph{clusters}. Two parameters describe the large-scale picture.
Infinite clusters exist for $p$ above the \emph{critical parameter}~$p_c$, and there is exactly one
infinite cluster for $p$ above the \emph{uniqueness threshold}~$p_u$. On the hypercubic lattice
$\Z^d$ there is never more than one infinite cluster~\cite{AKN87,BK89}, and the argument of Burton and
Keane extends this to every amenable transitive graph (see~\cite[Chapter~7]{LP16}); there the two
thresholds coincide. On the $3$-regular tree, in contrast, $p_c=1/2$ and $p_u=1$: for every
$p\in(1/2,1)$ there are infinitely many infinite clusters, kept apart by the branching geometry of the
tree. Benjamini and Schramm conjectured that the thresholds are separated on every quasi-transitive
graph that shares the decisive feature of the tree, namely that the boundary of every finite vertex
set is comparable to its size~\cite[Conjecture~6]{BS96}. Such graphs are called nonamenable.

These notes present a proof of this conjecture for transitive graphs, following the
manuscript~\cite{Ope26}. In informal terms, the main result reads as follows; the precise statement is
\cref{thm:main}.

\begin{quote}
\emph{Let $G$ be a connected, locally finite, simple, nonamenable transitive graph. Then the matrix of
connection probabilities at the critical parameter defines a bounded operator on $\ell^2(V)$.
Consequently $p_c<\ptwo\le p_u$, so there is a nonempty interval of parameters with infinitely many
infinite clusters, and percolation on $G$ has mean-field critical behaviour.}
\end{quote}

Here $\ptwo$ is the largest parameter up to which the connection probabilities define a bounded
operator on $\ell^2(V)$. This threshold was introduced by Hutchcroft~\cite{Hut19,Hut20L2}, who
conjectured that $p_c<\ptwo$ on every nonamenable quasi-transitive graph~\cite[Conjecture~1.3]{Hut20L2}.
The inequality $\ptwo\le p_u$ is elementary: a unique infinite cluster forces all connection
probabilities to be bounded away from zero, which is incompatible with boundedness on $\ell^2(V)$
(\cref{lem:p22-pu}). The strict inequality $p_c<\ptwo$ is the substance of the result. It follows from
boundedness at $p_c$ itself, because a bounded connection operator remains bounded after a small
increase of the parameter (\cref{lem:sprinkling}).

\paragraph{Earlier results.}
Several classes of graphs were known to satisfy $p_c<p_u$. Benjamini and Schramm proved it for
one-ended nonamenable transitive planar graphs~\cite{BS01}, following Lalley's work on Fuchsian
groups~\cite{Lal98}. Schonmann treated highly nonamenable graphs~\cite{Sch01}, and Nachmias and Peres
nonamenable Cayley graphs of large girth~\cite{NP12}. Gaboriau proved $p_c<p_u$ for unimodular
transitive graphs carrying nonconstant harmonic Dirichlet functions~\cite{Gab05}. For Cayley graphs of
a given nonamenable group, Pak and Smirnova-Nagnibeda~\cite{PSN00} and Thom~\cite{Tho15} showed that
\emph{some} generating set has $p_c<p_u$; the result presented here applies to \emph{every} finite
generating set (\cref{cor:cayley}). In the operator framework, Hutchcroft proved $p_c<\ptwo$ for
nonamenable Gromov hyperbolic graphs and for graphs with a quasi-transitive nonunimodular group of
automorphisms~\cite{Hut19}, building on~\cite{Hut20nonuni}. Choi and Seo extended this to Cayley graphs
of acylindrically hyperbolic groups~\cite{CS25}. The upper inequality $\ptwo\le p_u$ can be an
equality~\cite{HP24}.

Critical percolation on a nonamenable transitive graph has no infinite cluster: this was proved by
Benjamini, Lyons, Peres and Schramm in the unimodular case~\cite{BLPS99a} and by Hutchcroft for all
quasi-transitive graphs of exponential growth~\cite{Hut16}. We use it throughout, since it makes
critical clusters finite. The behaviour of infinite clusters for $p_c<p<p_u$, including the
simultaneous statements across parameters, was developed by H\"aggstr\"om, Peres and
Schonmann~\cite{HP99,Sch99,HPS99}.

\paragraph{Methods.}
The proof uses the classical toolbox of percolation: the Harris and van den Berg--Kesten
inequalities~\cite{Har60,vdBK85}, the tree-graph bounds of Aizenman and Newman~\cite{AN84}, sharpness of
the phase transition~\cite{AB87,Men86,DCT16}, Russo's formula, and the mass-transport
principle~\cite{BLPS99b,AL07}. Two further tools enter in an essential way. The first is the
decision-tree covariance inequality of O'Donnell, Saks, Schramm and Servedio~\cite{OSSS05}, used as in
Hutchcroft's derivation of near-critical bounds from critical tails~\cite{Hut20exp}. The second is a
concentration estimate for exploration scores in the spirit of Hutchcroft's two-ghost
inequalities~\cite{Hut20loc,Hut21}. The genuinely new ingredients of~\cite{Ope26} are a construction
that splices independent critical clusters together (\cref{sec:maps}), whose combinatorics is related
to the finite modifications in Kozma's work on products of trees~\cite{Koz11}, and a deformation of the
connection kernel that penalises isolated open paths (\cref{sec:deformation}).

\paragraph{About these notes.}
The manuscript~\cite{Ope26} treats quasi-transitive graphs, possibly with loops and multiple edges.
We restrict to transitive simple graphs, which removes the averaging over vertex types, the reduction
to simple graphs by subdivision and several case distinctions, and lets the main line of the argument
stand out. The mathematics is that of~\cite{Ope26}; we have re-derived every step while rewriting it,
specialised the statements to the transitive case, and reorganised the material around four steps
whose precise statements are collected in \cref{sec:strategy}. Where a proof deviates from~\cite{Ope26}
beyond this specialisation, we say so. Classical results are cited rather than reproved; proofs are
given for the less standard variants we need, some of them in \cref{app:aux}.

\paragraph{Outline.}
\Cref{sec:setting} fixes the setting and states the main results:
\cref{thm:main}, its consequence for Cayley graphs, and the mean-field consequences in
\cref{thm:consequences}. \Cref{sec:strategy} is an overview of the whole proof. It reduces
\cref{thm:main} to the boundedness of the critical connection operator (\cref{thm:crit-bdd}), explains
the mechanism, states the four steps precisely, and proves \cref{thm:crit-bdd} from them.
\Cref{sec:prelim} collects the tools. \Cref{sec:bridges,sec:maps,sec:moments,sec:triangle} prove the a
priori estimates on critical and near-critical clusters (Steps~1 and~2), and
\cref{sec:deformation,sec:escapes} the properties of the deformation (Steps~3 and~4).
\Cref{sec:consequences} derives \cref{thm:consequences}.
